% Line CSVs have quoted headers, so columns are referenced by index: 0 arc_length, 1 accel_X, 2 effective_plastic_strain, 3 vel_X.
% Acceleration (left axis, solid) and velocity (right axis, dashed) overlaid with two axis environments.
% The legend lives on the first axis; velocity entries are added there with \addlegendimage.
% Run 0 is drawn thicker so it stays visible where run 1 overlaps it.
\pgfplotsset{
    ssbdual/.style={
            width=0.8\linewidth,
            height=0.65\linewidth,
            xmin=0, xmax=1,
            scaled y ticks=base 10:3,
        },
    ssbrun0/.style={blue, line width=1.6pt},
    ssbrun1/.style={red, line width=0.8pt},
}
\begin{figure}[htbp]
    \centering
    \begin{subfigure}[t]{0.48\textwidth}
        \centering
        \begin{tikzpicture}
            \begin{axis}[
                    ssbdual,
                    axis y line*=left,
                    xlabel={Arc length},
                    ylabel={$a_x$},
                    grid=major,
                    legend pos=north east,
                    legend cell align=left,
                    legend style={font=\scriptsize},
                ]
                \addplot[ssbrun0] table[col sep=comma, x index=0, y index=1] {results/data/simpleshearband/run0_t0_line.csv};
                \addlegendentry{$a_x$, run 0}
                \addplot[ssbrun1] table[col sep=comma, x index=0, y index=1] {results/data/simpleshearband/run1_t0_line.csv};
                \addlegendentry{$a_x$, run 1}
                \addlegendimage{ssbrun0, dashed}\addlegendentry{$v_x$, run 0}
                \addlegendimage{ssbrun1, dashed}\addlegendentry{$v_x$, run 1}
            \end{axis}
            \begin{axis}[
                    ssbdual,
                    axis y line*=right,
                    axis x line=none,
                    ylabel={$v_x$},
                ]
                \addplot[ssbrun0, dashed] table[col sep=comma, x index=0, y index=3] {results/data/simpleshearband/run0_t0_line.csv};
                \addplot[ssbrun1, dashed] table[col sep=comma, x index=0, y index=3] {results/data/simpleshearband/run1_t0_line.csv};
            \end{axis}
        \end{tikzpicture}
        \caption{Restart step}\label{fig:ssb_av_t0}
    \end{subfigure}
    \hfill
    \begin{subfigure}[t]{0.48\textwidth}
        \centering
        \begin{tikzpicture}
            \begin{axis}[
                    ssbdual,
                    axis y line*=left,
                    xlabel={Arc length},
                    ylabel={$a_x$},
                    grid=major,
                    legend pos=north east,
                    legend cell align=left,
                    legend style={font=\scriptsize},
                ]
                \addplot[ssbrun0] table[col sep=comma, x index=0, y index=1] {results/data/simpleshearband/run0_t1_line.csv};
                \addlegendentry{$a_x$, run 0}
                \addplot[ssbrun1] table[col sep=comma, x index=0, y index=1] {results/data/simpleshearband/run1_t1_line.csv};
                \addlegendentry{$a_x$, run 1}
                \addlegendimage{ssbrun0, dashed}\addlegendentry{$v_x$, run 0}
                \addlegendimage{ssbrun1, dashed}\addlegendentry{$v_x$, run 1}
            \end{axis}
            \begin{axis}[
                    ssbdual,
                    axis y line*=right,
                    axis x line=none,
                    ylabel={$v_x$},
                ]
                \addplot[ssbrun0, dashed] table[col sep=comma, x index=0, y index=3] {results/data/simpleshearband/run0_t1_line.csv};
                \addplot[ssbrun1, dashed] table[col sep=comma, x index=0, y index=3] {results/data/simpleshearband/run1_t1_line.csv};
            \end{axis}
        \end{tikzpicture}
        \caption{Step after restart}\label{fig:ssb_av_t1}
    \end{subfigure}
    \caption{Acceleration (solid, left axis) and velocity (dashed, right axis) along the sample line for run 0 and run 1.}\label{fig:ssb_accel_vel}
\end{figure}
